% Generated draft: requires mathematical and source review.
\documentclass[11pt]{article}
\usepackage[margin=25mm]{geometry}
\usepackage{fontspec}
\setmainfont{Latin Modern Roman}
\usepackage{amsmath,amssymb,mathtools,unicode-math}
\setmathfont{Latin Modern Math}
\usepackage{enumitem,needspace,xurl,hyperref}
\setlength{\parindent}{0pt}
\setlength{\parskip}{5pt}
\setlength{\emergencystretch}{3em}
\newcommand{\N}{\mathbb{N}}
\newcommand{\Z}{\mathbb{Z}}
\newcommand{\Q}{\mathbb{Q}}
\newcommand{\R}{\mathbb{R}}
\newcommand{\C}{\mathbb{C}}
\newcommand{\F}{\mathbb{F}}
\newcommand{\A}{\mathbb{A}}
\newcommand{\PP}{\mathbb{P}}
\newcommand{\E}{\mathbb{E}}
\newcommand{\eps}{\varepsilon}
\newcommand{\dd}{\,\mathrm{d}}
\begin{document}
\section*{UnsolvedMath ID: 30005682}
\textbf{Source problem code:} catalogue entry
\subsection{Definitions and mathematical statement}
Let \(\N=\{1,2,3,\ldots\}\). A function \(f:\N\to\R\) is bounded if there exists \(B\in[0,\infty)\) such that \(|f(k)|\le B\) for every \(k\in\N\), and completely multiplicative if
\[
f(ab)=f(a)f(b)\qquad\text{for all }a,b\in\N.
\]
No additional normalization of \(f(1)\) is specified in the supplied statement.

Let \(P\in\Z[x,y]\) be homogeneous: there is an integer \(d\ge0\) such that every monomial of \(P\) with nonzero coefficient has total degree \(d\). Assume
\[
P(m,n)>0\qquad\text{for every }(m,n)\in\N^2.
\]
Thus every value \(P(m,n)\) appearing below belongs to the domain of \(f\). Positivity is required on positive integer pairs; no condition on the coordinate axes is stated. The source specifies no restriction on the degree beyond homogeneity.

\textbf{Question.} For every bounded completely multiplicative function \(f:\N\to\R\) and every homogeneous polynomial \(P\in\Z[x,y]\) satisfying the stated positivity condition, does there exist \(L=L(f,P)\in\R\) such that
\[
\lim_{\substack{N\to\infty\\N\in\N}}
\frac{1}{N^2}\sum_{m=1}^{N}\sum_{n=1}^{N}f\bigl(P(m,n)\bigr)=L?
\]

\textit{Source qualifications.} This formulation follows the supplied clean statement. The registry labels it \texttt{corrected\_verified} and reports reconstruction and checking against the official Oberwolfach report; that verification is supplied provenance, not an independent verification here. The original excerpt has corrupted mathematical notation and does not itself legibly supply all the hypotheses above. It also discusses a separate convergence theorem for multiplicative actions of Gaussian integers on a compact metric space, assuming finitely many transformations indexed by Gaussian primes and a unique common invariant probability measure, with convergence asserted for continuous observables at every point. That contextual theorem is not an additional hypothesis of this question. Its precise Gaussian domain and notation cannot be recovered reliably from the supplied excerpt. The excerpt expresses the expectation that new ideas will be needed.
\subsection{Short English statement}
For every bounded real completely multiplicative function \(f\) and every homogeneous integer polynomial \(P\) that is positive on positive integer pairs, must the average of \(f(P(m,n))\) over \(1\le m,n\le N\) converge as \(N\to\infty\)? The supplied record labels this question open; this draft does not independently verify that status.
\subsection{Sources}
\begin{itemize}
\item Local registry: \texttt{data/unsolvedmath/problems.json}, permanent UnsolvedMath ID \(30005682\), source problem code \texttt{catalogue entry}, title \emph{Convergence of Multiplicative Polynomial Averages}; supplied complete source record in the Oberwolfach Reports Open Problems collection. No individual proposer is identified. The supplied proposed year and year tag are 2023.
\item Supplied report citation: \emph{Aspects of Aperiodic Order (2024)}, Oberwolfach Reports, pp.\ 2107--2178, DOI \texttt{10.4171/owr/2023/37}. \url{https://doi.org/10.4171/owr/2023/37}. The record describes Oberwolfach Reports as documenting mathematical workshops and their open-problem discussions. The supplied citation uses 2024, while the DOI and registry year information contain 2023; their relationship has not been independently checked.
\item Supplied literature source, labeled ``Donoso-Le-Moreira-Sun'': \url{https://arxiv.org/abs/2309.07249}. No further bibliographic details are supplied.
\item Supplied literature assessment, checked on \texttt{2026-08-22}: the record labels the question open, describes the displayed average as a special case of the Frantzikinakis--Host question, reports that the Oberwolfach report calls that question still out of reach, and states that existing Gaussian-integer averages treat narrower multiplicative situations. The supplied checked-source list consists of the Oberwolfach report and the Donoso-Le-Moreira-Sun link above. This is supplied dated curation derived from the OWR record, not an independently verified literature assessment; time-sensitive status and cited claims require independent verification.
\end{itemize}
\end{document}
